\chapter{Prefazione}

\introduction{Mattia Brescia}{
    \textit{Il corso fornisce il linguaggio e gli strumenti di base della matematica per l'informatica:
    logica, teoria degli insiemi, insiemi ordinati, strutture algebriche, combinatoria, aritmetica, polinomi e grafi.}
  }
\pagestyle{plain}

\section{Come leggere questa dispensa}

\textbf{L'ordine degli argomenti.} Gli argomenti non seguono necessariamente l'ordine delle lezioni, ma un
\textbf{ordine topologico}: ogni definizione, teorema o proposizione compare \textbf{dopo} tutto ciò da cui dipende.
Leggendo dall'inizio alla fine non si incontra mai qualcosa che usa un concetto non ancora introdotto.

\textbf{La numerazione.} Definizioni, teoremi, proposizioni, lemmi e corollari condividono un'unica numerazione per
capitolo (ad esempio Definizione 1.3, Proposizione 1.4, Teorema 1.5), così è facile ritrovarli.

\textbf{Il box «Prerequisiti».} All'inizio di ogni teorema (o proposizione, lemma, corollario) c'è un elenco di
\textbf{cosa serve sapere} per capirlo al meglio: definizioni, altri risultati o sezioni. Ogni voce è
\textbf{cliccabile} e riporta la pagina in cui si trova; se una voce non ti è chiara, conviene ripassarla prima di
leggere la dimostrazione.

\begin{concetto}
\begin{Proposizione}{Esempio di enunciato}{esempio-legenda}
    \begin{prerequisiti}
        \dipvoce{Qui trovi l'elenco di cosa serve sapere, con numero, titolo e pagina di ogni voce.}
    \end{prerequisiti}
    Qui trovi l'enunciato.
\end{Proposizione}
\begin{dimostrazione}
    Qui trovi la dimostrazione, che termina con il simbolo a destra.
\end{dimostrazione}
\end{concetto}

\textbf{I box.}
\begin{itemize}
    \item \textcolor{red!75!black}{\textbf{Rosso}}: definizioni.
    \item \textcolor{primary!80!black}{\textbf{Blu}}: teoremi, proposizioni e lemmi, seguiti dalla dimostrazione.
    \item \textcolor{seagreen!65!black}{\textbf{Verde scuro}}: corollari.
    \item \textcolor{green-gradient-6!70!black}{\textbf{Verde chiaro}}: note e osservazioni.
    \item \textcolor{yellow-gradient-6!70!black}{\textbf{Giallo}}: attenzione, errori comuni.
    \item Gli esempi sono racchiusi tra due linee orizzontali.
\end{itemize}

\textbf{Le tavole di verità.} Nelle tavole i valori sono $\vero$ (vero) e $\falso$ (falso). Le colonne evidenziate
sono quelle da guardare: per una tautologia, la colonna deve contenere solo $\vero$; per un'equivalenza, le colonne
con lo \textbf{stesso colore} devono coincidere riga per riga.
